\section{Preliminaries}
In this section, we briefly describe preliminary background and the tools used in this paper.

\subsection{2-Prover 1-Round Games}

\begin{definition} \label{def:2p1r} A 2P1R Game ${\cal G}(\calV, \calU, \mu, {\cal R}, {\cal S},
\{ \pi_{VU}\} )$ consists of sets of questions $\calV, \calU$ and
sets of answers ${\cal R}, {\cal S}$ for the two provers
respectively, a distribution $\mu$ on the set of question pairs
$\calV \times \calU$ and for every question pair $(V,U)$ in the
support of $\mu$, a predicate $\pi_{VU}: {\cal R} \times {\cal S}
\mapsto \{0,1\}$ that defines the pairs of accepting answers.
A strategy of the %
provers is a map $\phi: \calV
\mapsto {\cal R}, \phi: \calU \mapsto {\cal S}$. The value of the
strategy $\phi$ is:
\[
  \val(\phi, {\cal G}) := \Pr_{(V,U) \sim \mu} \left[ \pi_{VU} (
  \phi(V), \phi(U) ) =1
 \right]\,.
\]
The value of the game ${\rm val}({\cal G})$ is the maximum value of
any prover strategy. A Projection Game is one where for every answer
of the first prover, there is exactly one accepting answer of the
second prover. For a Projection Game, the predicate $\pi_{VU}$ can
be thought of as a map $\pi_{VU}: {\cal R} \mapsto {\cal S}$ and
the accepting answers are of the form $(r, \pi_{VU}(r))$ for $r \in {\cal R}$.
For a Projection Game, $|{\cal S}| \leq |{\cal R}|$.
\end{definition}

A 2P1R Game is
best viewed as a game between the two provers and a verifier. The verifier
picks a random question pair $(V,U)$
from the distribution $\mu$, asks one question each to the two provers respectively,
and accepts if and only if the provers' answers satisfy the predicate $\pi_{VU}$.
The probability of acceptance
%
by
the verifier is same as the value of 
%
the
provers' strategy.

\begin{definition} Given a 2P1R Game ${\cal G}(\calV, \calU, \mu, {\cal R}, {\cal S},
\{ \pi_{VU}\} )$, the $k$-wise repeated game is
 \[
{\cal G}^{\otimes k}(\calV^k, \calU^k, \mu^k, {\cal R}^k, {\cal S}^k,
\{ \pi^k_{V'U'}\} )\,,
\] where for $V' = (V_1, \ldots, V_k)$ and $U' =
(U_1,\ldots, U_k)$, $\pi^k_{V'U'} := \bigwedge_{i=1}^k \pi_{V_iU_i}$.
\end{definition}

We state below Raz's Parallel Repetition Theorem along with the recent
improvements (and simplifications) by Holenstein and Rao.
\begin{theorem}[\cite{Raz, Holenstein, ARao}]\label{thm:pr} 
There exists an absolute constant $c > 0$ such that for a 2P1R
Game ${\cal G}$ with ${\rm val}({\cal G})= 1-\eps$,
 $$ {\rm val}({\cal G}^k) \leq (1-\eps^{3})^{c k}\,. $$
For a Projection Game, the bound of $(1-\eps^{2})^{c k}$ holds.
\end{theorem}

\subsection{Hardness of 3LIN}

Our reduction is from the 3LIN problem over a finite field. For the proof of \expref{Theorem}{thm:main}, we use H\aa stad's well-known hardness result for 3LIN~\cite{Has97}.
For the proof of Theorem~1.4, we need a hardness result for
3LIN with very good completeness and we use a result of Khot and Ponnuswami~\cite{KhotPonnuswami}.

\begin{definition} For a prime power $q$,
an instance $(X, \Phi)$ of 3LIN consists of a set of variables $X$
over $\F_q$ and a set of linear constraints $\Phi$ such that each
constraint depends on exactly three variables. Let $\OPT(X,
\Phi)$ denote the maximum fraction of constraints satisfied by any
assignment.
\end{definition}

\begin{theorem}[\cite{Has97}]\label{thm:Has97} 
For every constant $\eta > 0$ and a prime $q$, it is
\cclass{NP}-hard to distinguish whether a 3LIN instance $(X, \Phi)$ over
$\F_q$ has $\OPT(X,\Phi) \geq 1-\eta$ or $\OPT(X,\Phi)
\leq {1}/{q}+\eta$.
\end{theorem}

\begin{theorem}[\cite{KhotPonnuswami}]\label{thm:kp} 
There is a reduction from a SAT instance $\Gamma$
with size $n$ to a 3LIN instance $(X,\Phi)$ with size $N$ over a field $\F_q$ of characteristic $2$
such that: \begin{itemize}
\item $N$ and the running time of the reduction are bounded by $2^{O(\log^2 n)}$.
\item If $\Gamma$ is satisfiable, then $\displaystyle\OPT(X, \Phi) \geq
1- 2^{-\Omega(\sqrt{\log N})}$.
\item If $\Gamma$ is unsatisfiable, then $\displaystyle\OPT(X, \Phi) \leq
1- \Omega\left( \frac{1}{\log^3 N}  \right)$.
\end{itemize}
\end{theorem}

\begin{remark} \expref{Theorem}{thm:kp} is stated in~\cite{KhotPonnuswami} over the field $\F_2$. It also holds
for any extension field $\F_q$ since in the YES Case, the good assignment is over
$\F_2$ and in the NO Case, any assignment over $\F_q$ can be mapped
to an equally good assignment over $\F_2$, \eg, by taking the last bit of
every $\F_q$-element.
\end{remark}


\begin{remark} We present the proof of \expref{Theorem}{thm:main} by constructing a PCP over the prime field $\F_q$ and the Fourier analysis is over
this field. For the proof of Theorem~1.4, we use the same PCP but over
a field of characteristic $2$. We will omit the straightforward extension of our (Fourier) analysis to this case.
\end{remark}



\subsection{Hadamard code and Fourier analysis}

Let $q$ be a prime,   $\omega := e^{2\pi i/q}$ be the complex
$q^{th}$ root of unity and $\Omega := \{1, \omega, \ldots,
\omega^{q-1} \}$.

\begin{definition} 
%
%
The
%
Hadamard Code of $\alpha \in \F_q^n$ is defined
as the table of values of the linear function
 $\chi_\alpha: \F_q^n \mapsto \Omega$ where
  $$\forall x \in \F_q^n, \qquad  \chi_\alpha (x) = \omega^{\alpha \cdot  x}\,. $$
\end{definition}

The vector space of all functions $f : \F_q^n \mapsto {\mathbb C}$
has an orthonormal basis $\{ \chi_\alpha \mid \alpha \in \F_q^n\}$
where the inner product between two functions $f,g : \F_q^n \mapsto {\mathbb C}$ is defined as
(the expectation is over $x$ uniformly distributed in $\F_q^n$)
 $$ \left< f,g  \right> := \Exp_{x } \left[ f(x) \overline{g(x)}
 \right]\,. $$
Hence, every $f: \F_q^n \mapsto {\mathbb C}$ can be expressed
uniquely as
 $$ f = \sum_{\alpha \in \F_q^n}  \widehat{f}(\alpha) \ \chi_\alpha\,.
 $$
The coefficients $\widehat{f}(\alpha) \in {\mathbb C}$ are called
Fourier coefficients. These are defined by:
 $$ \widehat{f} (\alpha) = \left<f, \chi_\alpha \right> = \Exp_{x} \left[ f(x) \overline{\chi_\alpha(x)}
 \right]\,. $$
By Parseval's identity, $\sum_\alpha | \widehat{f}(\alpha)|^2
 = \| f\|_2^2 = \Exp_x [ |f(x)|^2]. $ In particular, for a function taking values in $\Omega$,
the sum of squared absolute values of all its Fourier coefficients equals $1$.


%
%
%
%


